% A RUNAWAY SHOWS WHAT IT HAD READ.
%
% When an \outer macro is met somewhere it may not stand, the reference
% first draws what it had read so far -- the kind on one line and the
% partial list on the next -- and only then the fault:
%
%   Runaway definition?
%   ->ab
%   ! Forbidden control sequence found while scanning definition of \zz.
%
% The kinds are `text' for a \message, \write, \special or \mark text,
% `definition' for a \def body, `argument' for a macro's argument and
% `preamble' for an alignment preamble. Each is shown from its own
% starting point: a definition from the `->' between what it matches and
% what it says, an argument from the brace that opened it, a text from
% nothing at all.
%
% The partial list is cut to within ten characters of the width the log
% prints to -- 69 of 79 -- and finished with \ETC.
%
% HSTeX drew the fault and nothing else.
%
% The offending name is put back to be read again, and what the scanner
% is holding becomes a SPACE, which joins whatever was being gathered --
% so a \message whose text ran away writes one space after what it had.
%
% The argument kind is not in this probe: the reference closes a runaway
% argument with an inserted \par and HSTeX with a }, which needs the
% whole macro call abandoned to match. See build/trip/ra2.tex and the
% comment at the insertion site.
\catcode`\{=1 \catcode`\}=2 \catcode`\#=6
\tracingonline=1
\outer\def\lz{Z}
\message{[A a runaway text]}
\message{XY\lz Q}
\message{[B a runaway definition]}
\def\zz{ab\lz cd}
\message{[C one with a parameter text]}
\def\ww#1{ab#1\lz cd}
\message{[D a long one, cut with ETC]}
\def\yy{XYZXYZXYZXYZXYZXYZXYZXYZXYZXYZXYZXYZXYZXYZXYZXYZXYZXYZXYZXYZXYZXYZXYZXYZXYZXYZ\lz Q}
\message{[E an empty one]}
\def\vv{\lz}
\message{[done]}
\end
